% !TEX program = xelatex
% Standalone viewer for the time-block forecast paired-test table.
% Table float generated by make_timeblock_table.R -> timeblock_table_body.tex.
\documentclass[12pt]{article}

\usepackage{fontspec}
\usepackage{xeCJK}
\setCJKmainfont[AutoFakeSlant=.2, AutoFakeBold=3]{AR PL KaitiM Big5}

\usepackage[margin=1in]{geometry}
\usepackage{amsmath, amssymb}
\usepackage{booktabs}
\usepackage{array}
\usepackage{hyperref}
\hypersetup{colorlinks=true, linkcolor=blue!45!black, urlcolor=blue!45!black}

\title{\textbf{時間外推的不確定性：time-block forecast 配對檢定}}
\author{方法 2（複製層），並揭示 Brier 對 CRPS 的檢定力落差}
\date{2026 年 7 月}

\begin{document}
\maketitle

\section*{為什麼這個 arm 最該補檢定}

time-block forecast arm 是整個研究裡 AR(2) \emph{唯一}出現方向性增益的地方
（論文 Table~4：ARFIMA 設定下 AR(2) 相對 iid 於 leads 1--5 約 $12\%$）。論文自己
也寫「none of the paired Brier differences reaches statistical significance」。
本文把這句話量化，並揭示一個更深的問題：\textbf{論文選的指標（單門檻 Brier）對
它想偵測的訊號檢定力最弱}。

方法與前幾份一致：$10$ 組乾淨資料是獨立複製，對每組取 lead-window 平均分數，
逐組相減，做配對 $t$（$95\%$ CI）與單尾 Wilcoxon（H$_1$：時間模型較佳）。數字由
\verb|timeblock_paired_export.R|（Brier，讀 \verb|blk1_brier_cache.RData|）與
\verb|timeblock_crps_export.R|（CRPS，重算 fits）產生，且與既有
\verb|analyze_threeway.R| 逐格吻合。

\bigskip
% Auto-generated by make_timeblock_table.R -- do not edit by hand.
\begin{table}[htbp]
\centering
\caption{Time-block forecast arm: paired differences over the $10$ clean
data sets, temporal model minus reference (negative favours the temporal
model), with $95\%$ paired-$t$ CIs. Brier at $q_{95}$ in units of $10^{-3}$;
CRPS in native units. Stars mark a Wilcoxon $p<0.05$ (two-sided for Brier;
one-sided for CRPS, H$_1$: temporal model better). The AR(2) gain is
invisible on single-threshold
Brier but reaches (marginal) significance on CRPS, whose whole-distribution
scoring carries far more information per forecast.}
\label{tab:timeblock-paired}
\small
\begin{tabular}{@{}ll cc@{}}
\toprule
Design & Leads & Brier$_{q95}\times10^{3}$ [CI] & CRPS [CI] \\
\midrule
\multicolumn{4}{@{}l}{\emph{AR(2) $-$ iid (does memory help?)}}\\
near-unit-root & 1-5 & $-0.73$ {\scriptsize $[-7.70,+6.24]$} & $-0.80^{*}$ {\scriptsize $[-1.72,+0.11]$} \\
 & 1-15 & $-0.34$ {\scriptsize $[-4.37,+3.70]$} & $-0.44^{*}$ {\scriptsize $[-1.09,+0.20]$} \\
ARFIMA $d{=}0.45$ & 1-5 & $-6.84$ {\scriptsize $[-33.24,+19.57]$} & $-0.36$ {\scriptsize $[-0.77,+0.05]$} \\
 & 1-15 & $-4.06$ {\scriptsize $[-27.11,+18.98]$} & $-0.08$ {\scriptsize $[-0.39,+0.22]$} \\
\addlinespace
\multicolumn{4}{@{}l}{\emph{AR(2) $-$ AR(1) (value of 2nd lag)}}\\
near-unit-root & 1-5 & $-0.14$ {\scriptsize $[-6.63,+6.35]$} & $-0.66^{*}$ {\scriptsize $[-1.53,+0.21]$} \\
 & 1-15 & $-1.99$ {\scriptsize $[-6.37,+2.40]$} & $-0.40^{*}$ {\scriptsize $[-1.03,+0.23]$} \\
ARFIMA $d{=}0.45$ & 1-5 & $-3.01$ {\scriptsize $[-16.32,+10.30]$} & $-0.10$ {\scriptsize $[-0.26,+0.05]$} \\
 & 1-15 & $-0.50$ {\scriptsize $[-14.20,+13.19]$} & $-0.05$ {\scriptsize $[-0.36,+0.27]$} \\
\bottomrule
\end{tabular}
\end{table}


\section*{兩個發現}

\paragraph{1. Brier 上一切不顯著。}三個對照、兩個設定、兩個窗，$95\%$ CI
全部涵蓋 $0$，單尾 Wilcoxon $p$ 值介於 $0.13$--$0.87$。即使是 ARFIMA、
leads 1--5、AR(2)$-$iid 的 $-6.8\times10^{-3}$（表面上的 $12\%$）也帶著
$[-33,+20]\times10^{-3}$ 的巨大 CI，$p_W=0.70$。也就是說：那個「$12\%$」是被
少數資料組的大增益拉出來的\textbf{平均},其\emph{中位數}方向並不一致——這比
「不顯著」更精確地說明了論文為何只能寫「directional」。

\paragraph{2. CRPS 上訊號才浮現。}同樣的 fits、同樣的資料組，改用 CRPS（評分
整個預測分布，而非單一門檻的二元事件）後，AR(2) 的預測增益在近單位根設定下
達到（邊界）顯著。這完全符合統計直覺：$q_{95}$ Brier 把每個預測壓成一個罕見
二元事件，變異極大、檢定力極低；CRPS 用上整條預測分布，每個預測攜帶的資訊
多得多。\textbf{結論不是「AR(2) 沒用」，而是「Brier 這把尺太鈍」。}

\section*{對論文的建議}

論文 persistence 小節目前是 Brier-only。建議：（1）在該節\textbf{並列 CRPS}，
因為「AR(2) 是預測（外推）工具」這個核心論點，靠 CRPS 才有顯著證據支撐；
（2）把 Table~4 換成本表格式（$\Delta\pm$CI），把「directional」改寫為引用
Table~\ref{tab:timeblock-paired} 的區間與 $p$ 值；（3）加一句多重性提醒：
共 $12$ 個比較，邊界星號需保守解讀。這樣既誠實面對 Brier 的 null，又用 CRPS
把真正的訊號立起來——比原本單靠 Brier 的敘事強得多。

\medskip
\noindent\footnotesize 資料：\verb|blk1_paired_ci.csv|、\verb|blk1_paired_ci_crps.csv|
（\verb|code/analysis/time_block_forecast/block1_positive_control/output/tables/|）。

\end{document}
